\DocumentMetadata{
  pdfversion=2.0,
  pdfstandard=ua-2,
  lang=en-US,
  tagging=on,
  tagging-setup={math/setup=mathml-SE,tabsorder=structure}
}
\documentclass[11pt,letterpaper]{article}
\usepackage[margin=0.68in,includefoot]{geometry}
\usepackage{newcomputermodern}
\usepackage{amsmath,amscd,mathtools}
\usepackage{tikz,adjustbox}
\providecommand{\square}{\mdlgwhtsquare}
\usepackage[hidelinks]{hyperref}
\hypersetup{
  pdftitle={Numerical Analysis PhD Exam, May 2000},
  pdfauthor={University of Florida Department of Mathematics},
  pdfsubject={Graduate mathematics examination},
  pdfdisplaydoctitle=true,
  bookmarksopen=true
}
\setcounter{secnumdepth}{0}
\setlength{\parindent}{0pt}
\setlength{\parskip}{0.28em}
\setlength{\emergencystretch}{3em}
\setlength{\leftmargini}{2.15em}
\setlength{\leftmarginii}{2.65em}
\setlength{\leftmarginiii}{3em}
\pagestyle{empty}
\raggedbottom
\allowdisplaybreaks
\NewTaggingSocketPlug{block/list/label}{examproblem}{%
  \tagstructbegin{tag=Lbl}\tagstructend%
  \tagstructbegin{tag=\LBody}%
  \tagstructbegin{tag=\ProblemHeadingTag,title-o={\ProblemHeadingText}}%
  \tagmcbegin{tag=\ProblemHeadingTag,actualtext={\ProblemHeadingText}}%
  #2%
  \tagmcend%
  \tagstructend%
}
\newenvironment{examproblems}[2]{%
  \def\ProblemHeadingTag{#2}%
  \AssignTaggingSocketPlug{block/list/label}{examproblem}%
  \begin{enumerate}%
  \setcounter{enumi}{#1}%
  \renewcommand{\labelenumi}{\textbf{\theenumi.}}%
  \setlength{\topsep}{0.35em}%
  \setlength{\partopsep}{0pt}%
  \setlength{\itemsep}{0.48em}%
  \setlength{\parsep}{0.18em}%
}{\end{enumerate}}
\newenvironment{examsubparts}{%
  \AssignTaggingSocketPlug{block/list/label}{default}%
  \begin{enumerate}%
  \setlength{\topsep}{0.18em}%
  \setlength{\partopsep}{0pt}%
  \setlength{\itemsep}{0.16em}%
  \setlength{\parsep}{0.1em}%
}{\end{enumerate}}
\newenvironment{examinstructions}{%
  \par\begingroup\itshape
}{%
  \par\endgroup\vspace{0.18em}
}
\makeatletter
\renewcommand\subsection{\@startsection{subsection}{2}{0pt}%
  {-1.25ex plus -.25ex minus -.1ex}%
  {0.55ex plus .1ex}%
  {\normalfont\large\bfseries}}
\renewcommand{\@maketitle}{%
  \newpage\null\vskip -1em%
  {\footnotesize\bfseries
    \href{https://www.ufl.edu/}{University of Florida}, \href{https://math.ufl.edu/}{Department of Mathematics}\par}%
  \vskip -0.5em%
  \tagstructbegin{tag=H1,title={Numerical Analysis PhD Exam, May 2000}}%
  \tagpdfparaOff%
  \tagmcbegin{tag=H1}%
    {\LARGE\bfseries\href{https://gma.math.ufl.edu/exams/numerical-analysis-phd/}{Numerical Analysis PhD Exam}, May 2000\par}%
  \tagmcend%
  \tagpdfparaOn%
  \tagstructend%
  \vskip 0.25em%
  \tagmcbegin{artifact}%
    \hrule height 1pt%
  \tagmcend%
  \vskip 1em%
}
\makeatother
\title{Numerical Analysis PhD Exam, May 2000}
\author{}
\date{}
\begin{document}
\maketitle
\thispagestyle{empty}
\begin{examproblems}{0}{H2}
\def\ProblemHeadingText{Problem 1.}
\item
This problem has two parts.
\begin{examsubparts}
\item[\textnormal{(a)}]
State the conditions under which a square matrix~\(A\) can be decomposed into a form \(A=SDS^{-1}\), where~\(D\) is diagonal.
\item[\textnormal{(b)}]
Can
\[
A=\begin{pmatrix}5&1\\0&3\end{pmatrix}
\]
 be diagonalized in the sense of (a)? Why?
\end{examsubparts}
\def\ProblemHeadingText{Problem 2.}
\item
This problem has three parts.
\begin{examsubparts}
\item[\textnormal{(a)}]
State the procedure for factoring a matrix~\(A\) into a form \(A=LU\), where~\(L\) is lower triangular, and~\(U\) is upper triangular.
\item[\textnormal{(b)}]
When can this type of decomposition be accomplished?
\item[\textnormal{(c)}]
What is the computational advantage of having lower and upper triangular matrices. An explanation without specific computation counts is sufficient.
\end{examsubparts}
\def\ProblemHeadingText{Problem 3.}
\item
Assume that a symmetric matrix can be factored into a Cholesky decomposition \(A=LL^t\), where~\(L\) is lower triangular.
\begin{examsubparts}
\item[\textnormal{(a)}]
Give an exact numerical method for finding the specific entries \(l_{i,j}\) of \(L=\{l_{i,j}\}\), where \(i=1:n\), and \(j=1:n\)?
\item[\textnormal{(b)}]
What numerical problems would you anticipate with this method?
\end{examsubparts}
\def\ProblemHeadingText{Problem 4.}
\item
This problem has three parts.
\begin{examsubparts}
\item[\textnormal{(a)}]
Describe the method for factoring a matrix~\(A\) into a form \(A=QR\), where~\(Q\) is orthogonal, and~\(R\) is upper triangular.
\item[\textnormal{(b)}]
When can such a decomposition be accomplished?
\item[\textnormal{(c)}]
What numerical “tricks” can be used to make it more stable?
\end{examsubparts}
\def\ProblemHeadingText{Problem 5.}
\item
This problem has three parts.
\begin{examsubparts}
\item[\textnormal{(a)}]
State the Gershgorin Circle Theorem.
\item[\textnormal{(b)}]
Prove the first assertion, i.e. that the eigenvalues must be within one of the circles.
\item[\textnormal{(c)}]
Let
\[
A=\begin{pmatrix}
5&.01&.03\\
.02&4&.01\\
.03&-.01&4
\end{pmatrix}.
\]
 What can you say, from the Gershgorin Circle Theorem, about the eigenvalues of~\(A\).
\end{examsubparts}
\def\ProblemHeadingText{Problem 6.}
\item
Suppose that a sequence is defined by \(x_{n+1}=g(x_n)\), where~\(g\) is continuous, and \(g([a,b])\subset[a,b]\).
\begin{examsubparts}
\item[\textnormal{(a)}]
Show that there is a fixed point \(\alpha\in[a,b]\) such that \(g(\alpha)=\alpha\).
\item[\textnormal{(b)}]
If in addition, \(|g'(x)|\leq k<1\) prove that this fixed point must be unique.
\item[\textnormal{(c)}]
Under the above assumptions show that the sequence must converge to the unique fixed point.
\end{examsubparts}
\def\ProblemHeadingText{Problem 7.}
\item
This problem has three parts.
\begin{examsubparts}
\item[\textnormal{(a)}]
Give an exact expression for a Lagrange polynomial of degree~\(n\), which would interpolate \(n+1\) given data points, \(f(x_k)\).
\item[\textnormal{(b)}]
Describe the differences between Hermite polynomials and Lagrange polynomials.
\item[\textnormal{(c)}]
Describe a procedure for using a limiting set of Lagrange polynomials, which interpolates a function at the points \(x_0,x_0+h,x_1,x_1+h,x_2,x_2+h,\ldots\) to construct the Hermite polynomial at the points \(x_0,x_1,x_2,\ldots\).
\end{examsubparts}
\def\ProblemHeadingText{Problem 8.}
\item
Given the Lagrange interpolation formula for \(f(x)\),
\[
f(x)=f(-h)\frac{x-h}{-2h}+f(h)\frac{x+h}{2h}+f''(\zeta(x))\frac{(x-h)(x+h)}{2!};
\]
\begin{examsubparts}
\item[\textnormal{(a)}]
Derive the trapezoid rule for \(\int_{-h}^{h}f(x)\,dx\),
\item[\textnormal{(b)}]
Derive the error formula for the trapezoid rule, and
\item[\textnormal{(c)}]
Describe the idea behind Gaussian Quadrature, and explain why it is more accurate than “traditional methods” such as trapezoid, Simpson’s, etc.
\end{examsubparts}
\def\ProblemHeadingText{Problem 9.}
\item
This problem has two parts.
\begin{examsubparts}
\item[\textnormal{(a)}]
Triple Recursion Formula. Let \(\{\phi_n(x)\}\) be an orthogonal family of polynomials on \([a,b]\), with weight function \(w(x)>0\). Then for \(n\geq1\), show that
\[
\phi_{n+1}(x)=(a_nx+b_n)\phi_n(x)-c_n\phi_{n-1}(x),
\]
 for some constants \(a_n,b_n\), and \(c_n\). Hint: Let \(g(x)=\phi_{n+1}(x)-A_nx\phi_n(x)\), where \(A_n\) is a constant chosen so that \(g(x)\) has degree~\(n\).
\item[\textnormal{(b)}]
What can you say about the location of the zeros of a family of orthogonal polynomials on \([a,b]\)?
\end{examsubparts}
\def\ProblemHeadingText{Problem 10.}
\item
This problem has two parts.
\begin{examsubparts}
\item[\textnormal{(a)}]
Describe Euler’s method for solving the initial value problem \(y'=f(t,y)\), \(y(a)=y_0\).
\item[\textnormal{(b)}]
Suppose that \(|f(t,x)-f(t,y)|\leq L|x-y|\), for a fixed constant~\(L\), and all~\(t\), \(x\) and~\(y\). Suppose further that all solutions \(y(t)\) to the initial value problem described in (a) satisfy \(|y''(t)|\leq M\), for some constant~\(M\). Then if \(w_i\) is the Euler solution to the initial value problem, prove that
\[
|y_{i+1}-w_{i+1}|\leq(1+hL)|y_i-w_i|+\frac{h^2M}{2}.
\]
\end{examsubparts}
\end{examproblems}
\end{document}
